% Slides from Lecture-8 lines 37--84 (main already inputs sections/cover before this file).
\begin{frame}[plain]
  \centering
  \includegraphics[width=0.8\paperwidth,height=0.75\paperheight,keepaspectratio]{assets/Lecture-8_Fine_Tuning_RLHF/slide_002.png}
\end{frame}

\begin{frame}[t]{Table of Contents}
  \scriptsize
  \begin{enumerate}
    \item \textcolor{red}{Motivation}
    \item Learning Outcomes
    \item \textcolor{red}{Fine-Tuning Methods for LLMs}
    \begin{enumerate}
      \item Categories of Fine-Tuning
      \item \textcolor{red}{Full Fine-Tuning (FT)}
      \item \textcolor{red}{Parameter-Efficient Fine-Tuning (PEFT)}
    \end{enumerate}
    \item \textcolor{red}{Supervised Fine-Tuning (SFT)}
    \begin{enumerate}
      \item What is SFT?
      \item Datasets for SFT
      \item Limitations of SFT
    \end{enumerate}
    \item \textcolor{red}{RLHF --- Reinforcement Learning with Human Feedback}
    \begin{enumerate}
      \item Why RLHF?
    \end{enumerate}
  \end{enumerate}
\end{frame}

\begin{frame}[t]{Table of Contents (Detailed Outline)}
  \scriptsize
  \begin{enumerate}
    \item[\textcolor{red}{2.}] RLHF Pipeline
    \item[\textcolor{red}{3.}] Reward Model (RM)
    \item[\textcolor{red}{4.}] RLHF Training (with PPO)
    \item[\textcolor{red}{5.}] Direct Preference Optimization (DPO)
    \item[\textcolor{red}{6.}] LoRA and Quantized LoRA
      \begin{enumerate}
        \item[1.] What is LoRA?
        \item[2.] Benefits of LoRA
        \item[3.] What is QLoRA?
        \item[4.] Applications of LoRA / QLoRA
      \end{enumerate}
    \item[\textcolor{red}{7.}] Limitations of Fine-Tuning and RLHF
    \item[\textcolor{red}{8.}] Future Directions
    \item[\textcolor{red}{9.}] Summary
  \end{enumerate}
\end{frame}
